% GENERATED FILE. Edit canonical lessons, then run npm run generate:books.
% canonical-source-hash: fnv1a64:985edf3a1bf12c92
% canonical-lessons: IT-C427-doganiere, IT-C427-dichiarare, IT-C427-perquisire, IT-C427-perquisizione, IT-C427-valico

\chapter{Al confine}
\label{ch:it-al-confine}
\hlchaptermodality{427}

\begin{chapteropening}
\textbf{By the end of this chapter:} \emph{I can get through border and security checks: the customs officer, declaring, being searched and the border pass.}

Complete the last lesson of chapter 427: I can get through border and security checks: the customs officer, declaring, being searched and the border pass.
\end{chapteropening}

\section[il doganiere]{il doganiere — the customs officer}

\label{lesson:IT-C427-doganiere}



\begin{quote}
Before the new one: say the Italian for the scales, then the Italian for to pack.
\end{quote}

\subsection*{You'll want to know: il doganiere}
\textbf{il doganiere} — \textquotedblleft{}the customs officer\textquotedblright{}.

\begin{quote}
Il doganiere apre la mia valigia.
\end{quote}

\begin{quote}
\emph{The customs officer opens my case.}
\end{quote}

\emph{La dogana} is customs.

\begin{cousinweb}
\emph{doganiere} is built on \emph{dogana} 'customs', which came from Arabic \emph{dīwān} 'register, office' through a southern form like Sicilian \emph{doana} (EtymDB).
\end{cousinweb}

\begin{grammarlens}[title={What you've built}]
One more word for telling, explaining and describing at B1.
\end{grammarlens}

\subsection*{Your turn}
\begin{itemize}
\raggedright
  \item \emph{Say it:} \emph{il doganiere}
  \item \emph{Say it:} \emph{il doganiere}, once more
  \item \emph{Say it:} a sentence of your own with \emph{il doganiere}
  \item \emph{Recall:} say the Italian for the scales, then the Italian for to pack, then say \emph{il doganiere} again
\end{itemize}

\begin{tcolorbox}[breakable,colback=teal!4,colframe=teal!35!black,title={Before you move on}]
What does il doganiere mean? (\textquotedblleft{}The customs officer\textquotedblright{}.) Say it once more.
\end{tcolorbox}
\section[dichiarare]{dichiarare — to declare}

\label{lesson:IT-C427-dichiarare}



\begin{quote}
Before the new one: say the Italian for to pack, then the Italian for the customs officer.
\end{quote}

\subsection*{You'll want to know: dichiarare}
\textbf{dichiarare} — \textquotedblleft{}to declare\textquotedblright{}.

\begin{quote}
Ha qualcosa da dichiarare?
\end{quote}

\begin{quote}
\emph{Do you have anything to declare?}
\end{quote}

\emph{Niente da dichiarare} — nothing to declare.

\begin{cousinweb}
\emph{dichiarare} is inherited from Latin \emph{dēclārāre} 'to make clear' (EtymDB), built on \emph{clārus} 'clear': Latin \emph{cl-} became \emph{chi-}, as in \emph{chiaro}.
\end{cousinweb}

\begin{grammarlens}[title={What you've built}]
One more word for telling, explaining and describing at B1.
\end{grammarlens}

\subsection*{Your turn}
\begin{itemize}
\raggedright
  \item \emph{Say it:} \emph{dichiarare}
  \item \emph{Say it:} \emph{dichiarare}, once more
  \item \emph{Say it:} a sentence of your own with \emph{dichiarare}
  \item \emph{Recall:} say the Italian for to pack, then the Italian for the customs officer, then say \emph{dichiarare} again
\end{itemize}

\begin{tcolorbox}[breakable,colback=teal!4,colframe=teal!35!black,title={Before you move on}]
What does dichiarare mean? (\textquotedblleft{}To declare\textquotedblright{}.) Say it once more.
\end{tcolorbox}
\section[perquisire]{perquisire — to search (a person or a place)}

\label{lesson:IT-C427-perquisire}



\begin{quote}
Before the new one: say the Italian for the customs officer, then the Italian for to declare.
\end{quote}

\subsection*{You'll want to know: perquisire}
\textbf{perquisire} — \textquotedblleft{}to search (a person or a place)\textquotedblright{}.

\begin{quote}
Al controllo mi hanno perquisito.
\end{quote}

\begin{quote}
\emph{At security they searched me.}
\end{quote}

Conjugated like \emph{finire}: \emph{perquisisco}.

\begin{cousinweb}
\emph{perquisire} is a learned borrowing from Latin \emph{perquīrere} 'to search everywhere', built on \emph{quaerere} 'to seek' (Lewis \& Short).
\end{cousinweb}

\begin{grammarlens}[title={What you've built}]
One more word for telling, explaining and describing at B1.
\end{grammarlens}

\subsection*{Your turn}
\begin{itemize}
\raggedright
  \item \emph{Say it:} \emph{perquisire}
  \item \emph{Say it:} \emph{perquisire}, once more
  \item \emph{Say it:} a sentence of your own with \emph{perquisire}
  \item \emph{Recall:} say the Italian for the customs officer, then the Italian for to declare, then say \emph{perquisire} again
\end{itemize}

\begin{tcolorbox}[breakable,colback=teal!4,colframe=teal!35!black,title={Before you move on}]
What does perquisire mean? (\textquotedblleft{}To search (a person or a place)\textquotedblright{}.) Say it once more.
\end{tcolorbox}
\section[la perquisizione]{la perquisizione — the search (by police or security)}

\label{lesson:IT-C427-perquisizione}



\begin{quote}
Before the new one: say the Italian for to declare, then the Italian for to search.
\end{quote}

\subsection*{You'll want to know: la perquisizione}
\textbf{la perquisizione} — \textquotedblleft{}the search (by police or security)\textquotedblright{}.

\begin{quote}
La perquisizione è durata dieci minuti.
\end{quote}

\begin{quote}
\emph{The search took ten minutes.}
\end{quote}

\emph{Un mandato di perquisizione} — a search warrant.

\begin{cousinweb}
\emph{perquisizione} is built on \emph{perquisire} 'to search', a learned borrowing from Latin \emph{perquīrere} 'to search everywhere', from \emph{quaerere} 'to seek' (Lewis \& Short).
\end{cousinweb}

\begin{grammarlens}[title={What you've built}]
One more word for telling, explaining and describing at B1.
\end{grammarlens}

\subsection*{Your turn}
\begin{itemize}
\raggedright
  \item \emph{Say it:} \emph{la perquisizione}
  \item \emph{Say it:} \emph{la perquisizione}, once more
  \item \emph{Say it:} a sentence of your own with \emph{la perquisizione}
  \item \emph{Recall:} say the Italian for to declare, then the Italian for to search, then say \emph{la perquisizione} again
\end{itemize}

\begin{tcolorbox}[breakable,colback=teal!4,colframe=teal!35!black,title={Before you move on}]
What does la perquisizione mean? (\textquotedblleft{}The search (by police or security)\textquotedblright{}.) Say it once more.
\end{tcolorbox}
\section[il valico]{il valico — the pass (in the mountains), the border crossing}

\label{lesson:IT-C427-valico}



\begin{quote}
Before the new one: say the Italian for to search, then the Italian for the search.
\end{quote}

\subsection*{You'll want to know: il valico}
\textbf{il valico} — \textquotedblleft{}the pass (in the mountains), the border crossing\textquotedblright{}.

\begin{quote}
Il valico è chiuso per la neve.
\end{quote}

\begin{quote}
\emph{The pass is closed because of the snow.}
\end{quote}

\emph{Il valico di frontiera} — the border crossing.

\begin{cousinweb}
\emph{valico} is built on \emph{valicare} 'to cross over', which is inherited from Latin \emph{vāricāre} 'to stand with legs apart, to straddle' (EtymDB).
\end{cousinweb}

\begin{grammarlens}[title={What you've built}]
One more word for telling, explaining and describing at B1.
\end{grammarlens}

\subsection*{Your turn}
\begin{itemize}
\raggedright
  \item \emph{Say it:} \emph{il valico}
  \item \emph{Say it:} \emph{il valico}, once more
  \item \emph{Say it:} a sentence of your own with \emph{il valico}
  \item \emph{Recall:} say the Italian for to search, then the Italian for the search, then say \emph{il valico} again
\end{itemize}

\begin{tcolorbox}[breakable,colback=teal!4,colframe=teal!35!black,title={Before you move on}]
What does il valico mean? (\textquotedblleft{}The pass (in the mountains), the border crossing\textquotedblright{}.) Say it once more.
\end{tcolorbox}
